\documentclass[11pt]{article}

\usepackage{amsmath,amssymb,amsthm,mathtools}
\usepackage[margin=1in]{geometry}
\usepackage{microtype}
\usepackage[hidelinks]{hyperref}

\newtheorem{theorem}{Theorem}[section]
\newtheorem{lemma}[theorem]{Lemma}
\newtheorem{corollary}[theorem]{Corollary}

\newcommand{\dd}{\,\mathrm{d}}
\newcommand{\one}{\mathbf{1}}
\newcommand{\ip}[2]{\left\langle #1,#2\right\rangle}
\newcommand{\join}{\mathbin{\vee}}

\title{Cones over the Paw and the Triangle--Edge Product}
\author{}
\date{}

\begin{document}

\maketitle

\begin{abstract}
Let $P$ be the paw graph, and let $R=K_3\sqcup K_2$.  For either
$B=P\sqcup\overline{K_s}$ or
$B=R\sqcup\overline{K_s}$, we prove that the cone
$H=K_1\join B$ satisfies
\[
t(H,W)\geq\Phi_H(p)
\]
for every graphon of edge density $p\geq2/3$.  This is the complete
required interval because $\chi(H)=4$.  The proof uses only the
previously proved paw and triangle bounds, multiplicativity, and a
fully proved weighted link argument.  It is not a general
universal-vertex closure.  The result proves the previously open Atlas
graphs 49, 156, and 165.
\end{abstract}


%%%%%%%%%%%%%%%%%%%%%%%%%%%%%%%%%%%%%%%%%%%%%%%%%%%%%%%%%%%%%%%%%%%%%%%%%%%%%%%
\section{The two base families and their common target}
%%%%%%%%%%%%%%%%%%%%%%%%%%%%%%%%%%%%%%%%%%%%%%%%%%%%%%%%%%%%%%%%%%%%%%%%%%%%%%%

Let $W\colon\Omega^2\to[0,1]$ be a symmetric measurable graphon on a
probability space $(\Omega,\mu)$.  For a finite graph $G$, write
\[
t(G,W)=
\int_{\Omega^{V(G)}}
\prod_{uv\in E(G)}W(x_u,x_v)
\prod_{u\in V(G)}\dd\mu(x_u).
\]
Put
\[
p=t(K_2,W),
\qquad
d(x)=\int_\Omega W(x,y)\dd\mu(y).
\]

The paw $P$ is a triangle with one pendant edge.  Let
\[
R=K_3\sqcup K_2.
\]
For $s\geq0$, define
\[
B_s^P=P\sqcup\overline{K_s},
\qquad
B_s^R=R\sqcup\overline{K_s}.
\]
The isolated vertices have no effect on homomorphism densities.

The chromatic polynomials are
\begin{align*}
\chi_{B_s^P}(x)
&=x^{s+1}(x-1)^2(x-2),
&v(B_s^P)&=s+4,
\\
\chi_{B_s^R}(x)
&=x^{s+2}(x-1)^2(x-2),
&v(B_s^R)&=s+5.
\end{align*}
Consequently, both base families have the same target
\begin{equation}
\Phi_{B_s^P}(z)=\Phi_{B_s^R}(z)
=\phi(z):=z^2(2z-1).
\label{eq:base-target}
\end{equation}

We use the following accepted positive inputs.

\begin{theorem}[Base bounds]
\label{thm:base}
For every graphon $V$ of edge density $z\geq1/2$,
\[
t(B_s^P,V)\geq\phi(z),
\qquad
t(B_s^R,V)\geq\phi(z).
\]
\end{theorem}

\begin{proof}
The paw is the member $L_1$ of the rooted triangle--tree theorem, so
\[
t(P,V)\geq z^2(2z-1).
\]
Adding isolated vertices changes neither side.  For the other family,
multiplicativity over disjoint unions and the triangle Goodman bound
give
\[
t(R,V)=t(K_3,V)t(K_2,V)
\geq z(2z-1)z=z^2(2z-1).
\]
Again, isolates make no change.
\end{proof}

The polynomial $\phi$ satisfies
\[
\phi(1/2)=0,
\qquad
\phi'(z)=2z(3z-1)>0,
\qquad
\phi''(z)=12z-2>0
\]
on $[1/2,1]$.  It is therefore nonnegative, increasing, and convex on
that interval.

\begin{lemma}[Global tangent bound for the bases]
\label{lem:tangent}
Let $B$ be any $B_s^P$ or $B_s^R$, and fix $c\in[1/2,1]$.  Every
graphon $V$, of arbitrary edge density $z\in[0,1]$, satisfies
\[
t(B,V)\geq\phi(c)+\phi'(c)(z-c).
\]
\end{lemma}

\begin{proof}
If $z\geq1/2$, Theorem~\ref{thm:base} and convexity give the result.
If $z<1/2$, the tangent line at $c$ is at most
$\phi(1/2)=0$ when evaluated at $1/2$, and it decreases as its argument
moves below $1/2$.  Thus the tangent is nonpositive at $z$, while
$t(B,V)\geq0$.
\end{proof}


%%%%%%%%%%%%%%%%%%%%%%%%%%%%%%%%%%%%%%%%%%%%%%%%%%%%%%%%%%%%%%%%%%%%%%%%%%%%%%%
\section{The weighted link estimate}
%%%%%%%%%%%%%%%%%%%%%%%%%%%%%%%%%%%%%%%%%%%%%%%%%%%%%%%%%%%%%%%%%%%%%%%%%%%%%%%

Define
\[
\tau(x)=
\int_{\Omega^2}W(x,y)W(x,z)W(y,z)
\dd\mu(y)\dd\mu(z).
\]

\begin{lemma}[Weighted rooted-triangle inequality]
\label{lem:weighted}
If $p>0$, then for every integer $h\geq0$,
\[
\int_\Omega d(x)^h\tau(x)\dd\mu(x)
\geq
\frac{2p-1}{p}
\int_\Omega d(x)^{h+2}\dd\mu(x).
\]
\end{lemma}

\begin{proof}
Put $M_j=\int d^j$, let $T_W$ be the graphon operator, and define
\[
I_h=\int d^h\tau,
\qquad
J_h=\ip{d^h}{T_Wd}.
\]
The inequality $ab\geq a+b-1$ on $[0,1]^2$, applied to
$W(x,y),W(x,z)$ and multiplied by $d(x)^hW(y,z)$, gives
\begin{equation}
I_h\geq2J_h-pM_h.
\label{eq:I-J}
\end{equation}

Set $U=1-W$, $u=T_U\one=1-d$, and $q=1-p$.  Then
\[
T_Wd=d-q\one+T_Uu,
\]
and hence
\begin{equation}
J_h=M_{h+1}-qM_h+\ip{d^h}{T_Uu}.
\label{eq:J}
\end{equation}
By symmetry of $U$, twice
\[
\ip{d^h}{T_Uu}-\int d^hu^2
\]
equals
\[
\int_{\Omega^2}U(x,y)
\bigl(u(y)-u(x)\bigr)
\bigl(d(x)^h-d(y)^h\bigr)
\dd\mu(x)\dd\mu(y),
\]
which is nonnegative because $d=1-u$.  Thus
\[
J_h\geq M_{h+2}-M_{h+1}+pM_h.
\]
Substitution into \eqref{eq:I-J} yields
\[
I_h\geq2M_{h+2}-2M_{h+1}+pM_h.
\]
Finally, subtracting $(2p-1)M_{h+2}/p$ from the right-hand side gives
\[
\frac1p\int_\Omega d(x)^h(d(x)-p)^2\dd\mu(x)\geq0.
\]
\end{proof}

For $d(x)>0$, set
\[
\dd\mu_x(y)=\frac{W(x,y)}{d(x)}\dd\mu(y),
\]
and let $W_x$ be the kernel $W$ on $(\Omega,\mu_x)$.  Its edge density
is $z_x=\tau(x)/d(x)^2$.  At $d(x)=0$, set $z_x=0$.

\begin{corollary}[Average density in the links]
\label{cor:average}
For every integer $N\geq2$,
\[
\int_\Omega d(x)^Nz_x\dd\mu(x)
\geq
\left(2-\frac1p\right)\int_\Omega d(x)^N\dd\mu(x).
\]
\end{corollary}

\begin{proof}
Apply Lemma~\ref{lem:weighted} with $h=N-2$.
\end{proof}


%%%%%%%%%%%%%%%%%%%%%%%%%%%%%%%%%%%%%%%%%%%%%%%%%%%%%%%%%%%%%%%%%%%%%%%%%%%%%%%
\section{The cone theorem}
%%%%%%%%%%%%%%%%%%%%%%%%%%%%%%%%%%%%%%%%%%%%%%%%%%%%%%%%%%%%%%%%%%%%%%%%%%%%%%%

\begin{theorem}[Special cones]
\label{thm:main}
Let $B$ be any $B_s^P$ or $B_s^R$, put $N=v(B)$, and define
\[
H=K_1\join B.
\]
Every graphon $W$ of edge density $p\geq2/3$ satisfies
\[
\boxed{t(H,W)\geq\Phi_H(p).}
\]
Moreover, the exact target is
\begin{equation}
\Phi_H(p)
=p^{N-3}(2p-1)^2(3p-2).
\label{eq:cone-target}
\end{equation}
Since $\chi(H)=4$, this is the complete required interval.
\end{theorem}

\begin{proof}
Put
\[
c=2-\frac1p.
\]
The hypothesis $p\geq2/3$ gives $c\in[1/2,1]$.  Conditioning on the
image $x$ of the cone vertex gives the exact identity
\[
t(H,W)=\int_\Omega d(x)^N t(B,W_x)\dd\mu(x).
\]
Let $M_N=\int d^N$.  Lemma~\ref{lem:tangent} gives
\begin{align*}
t(H,W)
&\geq
M_N\phi(c)
\\
&\quad+
\phi'(c)
\left(\int d(x)^Nz_x\dd\mu(x)-cM_N\right).
\end{align*}
The second line is nonnegative by Corollary~\ref{cor:average}, since
$c=2-1/p$ and $\phi'(c)\geq0$.  Therefore
\[
t(H,W)\geq M_N\phi(c)\geq p^N\phi(c),
\]
where the last step is Jensen and $\phi(c)\geq0$.

Now
\[
c=\frac{2p-1}{p},
\qquad
2c-1=\frac{3p-2}{p}.
\]
Consequently,
\[
p^N\phi(c)
=p^Nc^2(2c-1)
=p^{N-3}(2p-1)^2(3p-2).
\]

It remains to verify that this is the chromatic target.  For a cone,
\[
\chi_H(x)=x\chi_B(x-1).
\]
If $q=1-p$, then $1/q-1=p/q$.  Using
\eqref{eq:base-target} in its defining chromatic-polynomial form,
\begin{align*}
\Phi_H(p)
&=q^{N+1}\frac1q\chi_B\!\left(\frac pq\right)
\\
&=p^N\Phi_B\!\left(2-\frac1p\right)
=p^N\phi(c).
\end{align*}
This proves both the inequality and \eqref{eq:cone-target}.
\end{proof}


%%%%%%%%%%%%%%%%%%%%%%%%%%%%%%%%%%%%%%%%%%%%%%%%%%%%%%%%%%%%%%%%%%%%%%%%%%%%%%%
\section{Endpoints, equality, and catalogue consequences}
%%%%%%%%%%%%%%%%%%%%%%%%%%%%%%%%%%%%%%%%%%%%%%%%%%%%%%%%%%%%%%%%%%%%%%%%%%%%%%%

At $p=2/3$, the factor $3p-2$ vanishes, so the target is zero.  The
balanced complete tripartite graphon has no $K_4$ and attains equality,
because every graph in the theorem contains a $K_4$.  At $p=1$, both
sides equal $1$.  More generally, every balanced complete $k$-partite
graphon at $p=1-1/k$, $k\geq3$, attains equality by the chromatic
polynomial identity.

\begin{corollary}[New positive Atlas cases]
The following previously open graphs are positive on the full interval
$p\geq2/3$:

\begin{center}
\begin{tabular}{ccll}
\hline
Atlas ID & graph6 & base after deleting cone vertex & target \\
\hline
49  & \texttt{Df\char123}  & $P$ & $p(2p-1)^2(3p-2)$ \\
156 & \texttt{E\char126 H\_} & $P\sqcup K_1$ & $p^2(2p-1)^2(3p-2)$ \\
165 & \texttt{E\char126 aG} & $K_3\sqcup K_2$ & $p^2(2p-1)^2(3p-2)$ \\
\hline
\end{tabular}
\end{center}
\end{corollary}


%%%%%%%%%%%%%%%%%%%%%%%%%%%%%%%%%%%%%%%%%%%%%%%%%%%%%%%%%%%%%%%%%%%%%%%%%%%%%%%
\section{Scope and hidden assumptions}
%%%%%%%%%%%%%%%%%%%%%%%%%%%%%%%%%%%%%%%%%%%%%%%%%%%%%%%%%%%%%%%%%%%%%%%%%%%%%%%

The proof uses the rooted triangle--tree theorem for the paw, the
triangle Goodman bound, and multiplicativity over disjoint unions.
Every other analytic step is proved here.

This is not a general universal-vertex theorem.  It applies only because
the two stated base families share the explicitly verified target
$\phi(z)=z^2(2z-1)$ and the required graphon lower bound.  No conclusion
is asserted for a different base graph.

The link tangent is valid even when an individual link has density below
$1/2$: its tangent lower bound is then nonpositive.  Thus the proof does
not silently assume pointwise link density, regularity, minimum degree,
injective copies, or finite-step graphons.  It applies directly to every
graphon and ordinary homomorphism densities.

%%%%%%%%%%%%%%%%%%%%%%%%%%%%%%%%%%%%%%%%%%%%%%%%%%%%%%%%%%%%%%%%%%%%%%%%%%%%%%%
\appendix
\section{What the Lean formalization actually proves}
%%%%%%%%%%%%%%%%%%%%%%%%%%%%%%%%%%%%%%%%%%%%%%%%%%%%%%%%%%%%%%%%%%%%%%%%%%%%%%%

Theorem~\ref{thm:main} is formalized for both base families; Atlas $49$, $156$
and $165$ are \texttt{verified}.  The development is
\texttt{lean/Taeyoung/Methods/PawCone/}, over the abstract
\texttt{Methods/ConeBound.lean}.  Two differences of substance.

\subsection*{Theorem~\ref{thm:base} is proved, not accepted}

Both base bounds are formalized rather than taken as inputs: the paw is $L_1$
of the rooted triangle--tree development, and $t(R,V)=t(K_3,V)t(K_2,V)$ uses
Goodman, both of which this project proves.  So no accepted graph-density
result enters.

\subsection*{Lemma~\ref{lem:tangent} is proved without convexity}

$\phi''\geq0$ is never verified, and no convexity theory is used.  The tangent
bound is instead the polynomial identity
\[
 \phi(w)-\phi(c)-\phi'(c)(w-c)=(w-c)^2\bigl(2w+4c-1\bigr),
\]
whose right-hand side is nonnegative as soon as $w,c\geq1/2$; the case
$z<1/2$ is unchanged.  Since both families share the target
\eqref{eq:base-target}, everything after the base bound is proved once against
an abstract base graph, and the two families are two instantiations of it.

\end{document}
